% from ega4-17-fr.tex
\oldpage[IV-4]{74}

restriction de $f_1$, est encore de présentation finie et est en outre fidèlement plat; de plus, si $g_2=g_1|V$, $g_2\circ f_2$ est de type fini (resp. de présentation finie), l'ouvert $f_1^{-1}(V)$ étant quasi-compact (resp. quasi-compact et quasi-séparé). On est alors ramené aux hypothèses de (11.3.16), d'où l'on conclut que $g_2$ est un morphisme de type fini (resp. de présentation finie).

Dans le cas (iii) la question est locale sur $Z$, donc on peut supposer $Z$ affine, et il en résulte alors que $X$ est quasi-compact, donc aussi $Y=f(X)$. Le cas (iii) est donc conséquence de (i).

Dans le cas (iv) (avec les hypothèses supplémentaires sur $f$ ou $g$), on peut aussi supposer $Z$ affine, donc $X$ et $Y$ quasi-compacts; on sait en outre déjà que $g$ est localement de présentation finie et quasi-compact (1.1.3), donc tout revient à voir que $g$ est quasi-séparé, et il suffit donc de montrer que cette propriété est vraie lorsque l'on suppose $f$ quasi-compact. Or, puisque $g\circ f$ est quasi-séparé, il en est de même de $f$ (1.2.2, (v)); comme $f$ est quasi-compact et localement de présentation finie, il est de présentation finie (1.6.1); il suffit alors de répéter le raisonnement du premier alinéa de la démonstration de (11.3.16).

\begin{remark}[17.7.6]
\label{IV.17.7.6}
Dans l'assertion relative à la propriété (iv), on ne peut supprimer l'hypothèse que $f$ est quasi-compact; sans quoi, cela impliquerait que tout morphisme $g:Y\to Z$ quasi-compact et localement de présentation finie serait de présentation finie, conclusion que l'on sait être erronée (1.6.4). En effet, on peut se borner au cas où $Z$ est affine, donc $Y$ quasi-compact; il y a par suite un recouvrement fini $(U_i)$ de $Y$ par des ouverts affines tels que les restrictions $g|U_i$ soient de présentation finie; il suffirait alors de prendre pour $X$ le préschéma \emph{somme} des $U_i$, pour $f:X\to Y$ le morphisme canonique, qui est évidemment fidèlement plat et localement de présentation finie (1.4.3); $g\circ f$ serait de présentation finie en vertu du choix des $U_i$ et de (1.6.5), d'où notre assertion.
\end{remark}

\begin{proposition}[17.7.7]
\label{IV.17.7.7}
Soient $f:Y\to S$ et $h:X\to S$ deux morphismes localement de présentation finie, $g:X\to Y$ un $S$-morphisme, $x$ un point de $X$, $y=g(x)$. Supposons que $g$ soit plat au point $x$. Alors, si $h$ est lisse (resp. non ramifié, resp. étale) au point $x$, $f$ est lisse (resp. non ramifié, resp. étale) au point $y$.
\end{proposition}

Posons $s=h(x)=f(y)$. Dire que $h$ est non ramifié au point $x$ (resp. que $f$ est non ramifié au point $y$) équivaut à dire que $h^{-1}(s)$ est étale sur $k(s)$ au point $x$ (resp. que $f^{-1}(s)$ est étale sur $k(s)$ au point $y$) (17.4.1 et 17.6.1); comme le morphisme $g_s:h^{-1}(s)\to f^{-1}(s)$ déduit de $g$ est plat au point $x$, on voit que l'on peut se borner à prouver la proposition lorsque $h$ est lisse ou étale au point $x$. En outre, comme $h$ est alors plat au point $x$ (17.5.1), $f$ est plat au point $y$, comme il résulte de (2.2.11, (iv)). Il revient donc au même (17.5.1) de dire que $f$ est lisse (resp. étale) au point $y$, ou que $f^{-1}(s)$ est lisse (resp. étale) sur $k(s)$ au point $x$. On est ainsi ramené au cas où $S=\Spec(k)$ est le spectre d'un corps.

(i) \emph{Cas des morphismes lisses}. Comme $g$ est un morphisme localement de présentation finie (1.4.3, (v)), il y a un voisinage ouvert $U$ de $x$ dans $X$ dans lequel $g$ est plat (11.3.1) et $h$ lisse. En outre $g(U)$ est un voisinage ouvert de $y$ dans $Y$ (2.4.6); remplaçant $X$ par $U$ et $Y$ par $g(U)$, on peut donc supposer que $g$ est \emph{fidèlement plat} et que $h$ est lisse, et on est ramené à prouver que $f$ est alors lisse. Si $k'$ est une extension algébriquement

